\chapter{From ordinary statements to mathematical structures}
A mathematical proof is a sequence of statements whose conclusion follows
from explicitly declared assumptions. Start with the sentence
\emph{If a spike occurred in an interval, the interval contains a spike}.
The implication is meaningful only after the interval and the membership
relation have been defined.

\begin{definition}
A \emph{relation} from $A$ to $B$ is a subset $R\subseteq A\times B$.
It is a \emph{function} $f:A\to B$ when every $a\in A$ occurs as
the first coordinate of exactly one pair $(a,b)\in R$.
\end{definition}
For example, a spike train of a unit is a subset of its observation
time interval. A classification of spike waveforms by cluster labels is
a function only when every input waveform receives exactly one label.
This does not imply that the label identifies a physiological neuron.

\entry{1}{De Morgan's laws}
{For propositions $P,Q$, $\neg(P\land Q)\Leftrightarrow
(\neg P)\lor(\neg Q)$; for sets,
$(A\cup B)^c=A^c\cap B^c$, relative to an ambient set.}
{Transform selection criteria and complements of detected event sets.}
\entry{2}{Quantifier duality}
{For any predicate $P$, $\neg(\forall x\,P(x))\Leftrightarrow
\exists x\,\neg P(x)$; likewise
$\neg(\exists x\,P(x))\Leftrightarrow\forall x\,\neg P(x)$.}
{Express universal reliability claims and their counterexamples.}
\entry{3}{Mathematical induction (principle)}
{If $P(0)$ holds and $P(n)\Rightarrow P(n+1)$ for $n\in\N$,
then $P(n)$ holds for all $n\in\N$.}
{Prove recurrences, finite sums, binning bounds, and iterative invariants.}
\entry{4}{Recursion theorem on the natural numbers}
{With a starting value and a step map, there is a unique
sequence satisfying the corresponding recursive definition
in the standard set-theoretic setting.}
{Construct discrete filters, numerical solvers, and event counters.}
\entry{5}{Equivalence relations and partitions}
{Reflexive, symmetric, transitive relations correspond to partitions
into disjoint nonempty equivalence classes.}
{Distinguish clustering as a partition from physiological identification.}
\entry{6}{Well-ordering principle}
{Every nonempty subset of $\N$ has a least element.}
{Prove termination of monotone integer-index algorithms.}
\entry{7}{Cantor--Schroeder--Bernstein theorem}
{If injections $A\to B$ and $B\to A$ exist, then a bijection
$A\to B$ exists.}
{Understand cardinality and information-preserving relabelings.}
\entry{8}{Cantor's diagonal argument}
{There is no surjection $\N\to(0,1)$; the real interval
is uncountable.}
{Understand the distinction between finite sampled arrays and continuum models.}
\entry{9}{Least-upper-bound property (completeness axiom for $\R$)}
{Every nonempty subset of $\R$ bounded above has a real supremum.}
{Build limits, continuity, completeness, and convergence of numerical approximations.}
\entry{10}{Archimedean property and density of $\Q$}
{For real $x$ there is an integer $n>x$; between any
$a<b$ in $\R$ lies some $q\in\Q$.}
{Approximate continuous times and amplitudes by rational values, without confusing
approximation with exact sampling.}

\begin{example}[Equivalence classes of spike timestamps]
For a fixed bin width $\Delta>0$, let $s\sim t$ when
$\lfloor s/\Delta\rfloor=\lfloor t/\Delta\rfloor$.
This relation is reflexive, symmetric, and transitive.
Each equivalence class represents one temporal bin, not one spike.
A spike-count sequence counts members of each class.
\end{example}

\begin{caution}
Membership ($x\in A$), inclusion ($A\subseteq B$), and association
under a binary relation ($xRy$) are distinct concepts. An observed
cross-correlation is not evidence that a causal or anatomical edge exists.
\end{caution}
